% Complete reconstruction of the assigned original Theorem 6 dependencies.

\begin{annexitem}{annex:p9.n-c-004.1.theorem6}{thm:p09-fagin74-theorem6}{Both directions of the original generalized-spectrum theorem}
\begin{proof}
Throughout the input type \(S\) is fixed and relational, its relations have
positive arity, and every universe is finite and nonempty.  Write its size
as \(n\).  The argument follows Fagin's proof on printed pp.~53--58
\cite{Fagin1974}, with the standard reductions and the tableau's typing and
carry conditions made explicit.

\paragraph{Auxiliary functions and prenex form.}
Replace each auxiliary \(d\)-ary function \(f\) by a fresh
\((d+1)\)-ary relation \(F_f\) and require
\(\forall\bar x\exists!y\,F_f(\bar x,y)\).  For a constant use a unary
relation with exactly one element.  Define a term-value formula recursively:
\(V_x(z)\) is \(z=x\), a constant's value is its singleton predicate,
and
\[
 V_{f(t_1,\ldots,t_d)}(z):=
 \exists z_1\cdots z_d
  \left(\bigwedge_j V_{t_j}(z_j)\land F_f(z_1,\ldots,z_d,z)\right).
\]
Use fresh variables at each occurrence.  The graph axioms give a unique term
value by induction on the term.  Replace atomic formulas by the formula
which selects those values and applies the original relation or equality;
extend through Boolean operations and first-order quantifiers.  Formula
induction proves satisfaction preservation.  Conversely every graph expansion
defines exactly the functions of a function expansion on the same universe.
Thus this conversion preserves the class of reducts in both directions.

Rename bound variables before moving quantifiers to the front.  The steps
\((Qx\alpha)\land\beta\longleftrightarrow Qx(\alpha\land\beta)\)
and the analogous disjunction steps apply when \(x\) is not free in
\(\beta\); negation interchanges \(\forall\) and \(\exists\).
Structural induction gives a fixed prenex sentence
\(Q_1x_1\cdots Q_mx_m\varphi\), with quantifier-free \(\varphi\).
If one common arity is wanted as in the source, pad a \(d\)-ary relation
to arity \(h\geq d\) and require independence of the last \(h-d\)
coordinates.  Nonemptiness supplies those coordinates; taking any such
coordinates recovers the original relation.  This is a two-way same-universe
coding, not a change of the input type.

\paragraph{Spectrum to machine: a single guess followed by deterministic evaluation.}
First parse the binary \(n\), separators and the required relation words.
Use counters to check the \(n^{m_i}\) lengths and reject if a word is not
an encoding.  Guess every bit of each auxiliary relation table, once, in
lexicographic order.  The number of bits is \(\sum_j n^{d_j}\) for fixed
arities \(d_j\).  Evaluate the quantifier-free matrix at each required
assignment by computing the lexicographic index of each atomic tuple and
reading its input or guessed bit.  Even a linear scan of an entire tape
per access takes time polynomial in \(n\).

Evaluate a universal quantifier by taking the conjunction of its \(n\)
child values and an existential quantifier by taking their disjunction.
Store the assignment stack and the partial conjunction/disjunction for each
level.  There are at most \(1+n+\cdots+n^m\) nodes.  The matrix size,
\(m\), table arities and number of tables are fixed, so the entire
verification takes \(n^{O(1)}\) time.  It accepts exactly the expansions
satisfying the sentence.  In particular the existential relation guess is
\emph{outside} the first-order assignment loop; a different table at each
assignment would verify a different formula.

If \(S\ne\varnothing\), the encoding length \(\ell\) is at least
\(n\), so this is polynomial in \(\ell\).  If \(S=\varnothing\),
\(n<2^\ell\), so \(n^d<2^{d\ell}\).  This is exactly the original
\(NP_1\) scale.  Constant factors and finitely many short inputs can be
absorbed by increasing the exponent.  Malformed inputs need not be searched
indefinitely: the parsing counters are explicitly bounded by the represented
\(n\) and fixed arities.

\paragraph{Preparing the converse without changing its input convention.}
An \(NP\) machine on the original nonempty-type encodings takes
\(\ell^{O(1)}=n^{O(1)}\) time.  An \(NP_1\) machine on binary size
inputs takes \(2^{O(\ell)}=n^{O(1)}\) time.  We may clock the machine
and reject when its polynomial-in-\(n\) budget expires; accepting inputs
retain a bounded accepting branch.  For the construction use the source's
modified \(r+1\)-tape input: tape~1 contains the binary representation
of \(n\) backwards, and tape \(i+1\) contains \(b_i\).
Reversal and concatenation, with separators, convert these tapes to the
original input in polynomial time in \(n\).

There is no need to add unboundedly many tapes to do this conversion.  On the
first tape shift the reversed numerical word rightwards to make space for a
work area; append the separate \(b_i\)'s there one bit at a time by scanning
their tapes; reverse the numerical prefix in the work area.  Then simulate
the original fixed-tape machine using a finite product alphabet whose tracks
hold all its tapes and head marks.  A sweep finds the scanned symbols and
marked heads; another sweep writes the chosen transition and moves the marks.
One simulated step uses at most a polynomial number of sweeps over the
already visited interval.  If a boundary needs extending, shift the finite
work word and insert a fresh blank symbol.  The interval length is at most
input length plus simulated time, both polynomial in \(n\).  The sweep
simulation therefore has a polynomial time bound.  For \(r=0\) this is a
one-tape simulation.  Its finite alphabet and state set are fixed.
The source's convention that a scanned square is not rewritten as the actual
blank can be met by a separate nonblank symbol coding a simulated blank.
These elementary simulations suffice here; no exact speed-up factor is used.

Call the resulting machine \(M\), its state set \(K\), and its tape
alphabet \(\Gamma\), containing distinct \(0,1,B\).
Its tapes are one-way infinite, heads start at the first square, and a
left move from the first square rejects.  Increase a fixed integer \(k\)
so that \(k>m_i\) for every input arity and, for all sufficiently large
\(n\), an accepting run, its initial words and all scanned squares fit
strictly within \(n^k\) cells and \(n^k-1\) transitions.  There are
only finitely many remaining sizes.  We treat them separately below.

\paragraph{Order and finite codes: the source's first three blocks.}
Existentially quantify a strict linear order \(<\), endpoints \(c_1,c_2\),
and a unary function \(S\).  The first-order matrix says
\[
 \begin{gathered}
 \forall x\neg(x<x),\quad
 \forall x\forall y\forall z((x<y\land y<z)\to x<z),\\
 \forall x\forall y(x=y\lor x<y\lor y<x),\quad
 \forall x\neg(x<c_1),\quad\forall x\neg(c_2<x),\\
 S(c_2)=c_1,\qquad
 \forall x(x\ne c_2\to
    (x<Sx\land\neg\exists z(x<z\land z<Sx))).
 \end{gathered}
\]
Finite nonempty linear orders have unique endpoints and successors, so these
axioms give the intended cyclic successor.  Choose pairwise distinct tape
symbol codes and pairwise distinct state codes from the first
\(\lvert\Gamma\rvert\) and \(\lvert K\rvert\) successive elements.
The two code lists may overlap: a tape symbol and a state are separate typed
uses.  This works when \(n\geq\max(\lvert\Gamma\rvert,\lvert K\rvert)\).
Order the tape-symbol list as \(0,1,B,\ldots\), and the state list as
\(q_0,q_A,q_R,\ldots\).  Thus their respective first three codes are
\(c_1,S(c_1),S^2(c_1)\), as in the source's first block.
Write \(\kappa_b\) for the fixed successor term coding state \(b\), and
\(\gamma_e\) for that coding tape symbol \(e\).  State names and tape
symbol names in subsequent formulas abbreviate these terms.
Explicit range constraints are
\[
 \forall\bar t\bigvee_{b\in K}q(\bar t)=\kappa_b,\qquad
 \bigwedge_i\forall\bar t\forall\bar x
                \bigvee_{e\in\Gamma}v_i(\bar t;\bar x)=\gamma_e,
 \qquad
 \forall x\forall y\,(G(x,y)=0\lor G(x,y)=1\lor G(x,y)=B).
\]
Every conjunction and disjunction here is over a fixed finite set.

For tuples \(\bar x,\bar y\in A^k\), define lexicographic order by
\[
 \bar x<_{\rm lex}\bar y
 \quad\longleftrightarrow\quad
 \bigvee_{j=1}^{k}
  \left(\bigwedge_{h<j}x_h=y_h\land x_j<y_j\right).
\]
This is one fixed first-order formula.  Let
\(\mathbf a=(c_1,\ldots,c_1)\),
\(\mathbf z=(c_2,\ldots,c_2)\).  Require the \(2k\)-ary relation
\(S_1\) to be precisely immediate lexicographic succession, together with
the pair \((\mathbf z,\mathbf a)\).  Immediate succession is the formula
\(\bar x<_{\rm lex}\bar y\land
 \neg\exists\bar v(\bar x<_{\rm lex}\bar v\land
                       \bar v<_{\rm lex}\bar y)\).
Thus the tuples enumerate exactly \(N=n^k\) times and squares.
For actual time transitions put
\(D(\bar t,\bar u):=S_1(\bar t,\bar u)\land\bar t\ne\mathbf z\);
there is no final-to-initial time transition.

\paragraph{The binary-digit block, with its carry condition explicit.}
Use a binary function \(G(x,y)\), with values in \(\{0,1,B\}\).
Under the order, \(x\) names the integer from \(1\) to \(n\), while
\(y\) names a bit position from least significant upwards.  Initialize
\[
 G(c_1,c_1)=1,\qquad \forall y(y\ne c_1\to G(c_1,y)=B).
\]
Put \(C(x,y):=\forall z(z<y\to G(x,z)=1)\).
For \(x\ne c_2\), require for every \(y\)
\[
 \begin{array}{ll}
 \neg C(x,y) &\to G(Sx,y)=G(x,y),\\
 C(x,y)\land G(x,y)=1 &\to G(Sx,y)=0,\\
 C(x,y)\land(G(x,y)=0\lor G(x,y)=B)&\to G(Sx,y)=1.
 \end{array}
\]
These exhaustive, disjoint cases define the next row from the preceding row.
Induction on the ordered \(x\)'s shows that it is the reversed binary word
for that integer, followed by blanks: lower bits equal to~1 become~0 until
the first~0 becomes~1, or the first blank becomes a new leading~1; all
higher positions are unchanged.  There are enough positions because
\(\lfloor\log_2n\rfloor+1\leq n\).  This is the intended binary
counter in the source's block~4, with the all-lower-bits carry test written
explicitly rather than relying on overlapping printed antecedents.

\paragraph{States, heads and input tapes: blocks five through seven.}
Quantify \(q:A^k\to A\), \(v_i:A^{2k}\to A\), and
\(H_i\subseteq A^{2k}\), for \(1\leq i\leq r+1\), subject to the
range restrictions above.  At each time require exactly one tuple
\(\bar x\) with \(H_i(\bar t;\bar x)\).
Set \(q(\mathbf a)=q_0\), \(q(\mathbf z)=q_A\), and
\(H_i(\mathbf a;\mathbf a)\) for every tape.
Initialize tape~1 by
\[
 \begin{split}
 v_1(\mathbf a;c_1,\ldots,c_1,y)&=G(c_2,y),\\
 \neg(x_1=c_1\land\cdots\land x_{k-1}=c_1)
    &\to v_1(\mathbf a;\bar x)=B.
 \end{split}
\]
This gives the reversed binary word for \(n\), then blanks.
For tape \(i+1\), set at positions
\(\bar x=(c_1,\ldots,c_1,x_1,\ldots,x_{m_i})\)
\[
 v_{i+1}(\mathbf a;\bar x)=1\ \longleftrightarrow\
 P_i(x_1,\ldots,x_{m_i}),
 \qquad
 v_{i+1}(\mathbf a;\bar x)=0\ \longleftrightarrow\
 \neg P_i(x_1,\ldots,x_{m_i}).
\]
At all positions outside this initial segment require blank.  These are
first-order formulas and list exactly \(b_i\) in tuple order.
If \(r=0\), these relation-tape clauses are an empty conjunction;
the numerical tape and all other blocks remain.

\paragraph{Freezing final states and selecting a legal move: blocks eight and nine.}
For \(D(\bar t,\bar u)\) and \(q(\bar t)\in\{q_A,q_R\}\), require
\[
 q(\bar u)=q(\bar t),\quad
 \forall\bar x\,v_i(\bar u;\bar x)=v_i(\bar t;\bar x),\quad
 \forall\bar x(H_i(\bar u;\bar x)\longleftrightarrow
                              H_i(\bar t;\bar x))
\]
for every tape.  For each nonfinal state \(b\) and scanned-symbol list
\((e_1,\ldots,e_{r+1})\in\Gamma^{r+1}\), require
\[
 \begin{split}
 D(\bar t,\bar u)\land q(\bar t)=b\land
 \bigwedge_i\bigl(H_i(\bar t;\bar x^i)\land
                    v_i(\bar t;\bar x^i)=e_i\bigr)
 \ \to\ \bigvee_{s\in\delta(b,e_1,\ldots,e_{r+1})}\Phi_s.
 \end{split}
\]
There are only finitely many such table entries, so take their conjunction.
For \(s=(b';d_1,\ldots,d_{r+1};T_1,\ldots,T_{r+1})\), its
first-order formula \(\Phi_s\) says:
\[
 \begin{gathered}
 q(\bar u)=b',\qquad
 \bigwedge_i v_i(\bar u;\bar x^i)=d_i,\\
 \bigwedge_i\forall\bar y\bigl(\bar y\ne\bar x^i\to
              v_i(\bar u;\bar y)=v_i(\bar t;\bar y)\bigr),\\
 \bigwedge_{i:T_i=R}\exists\bar y
     \bigl(\bar x^i\ne\mathbf z\land S_1(\bar x^i,\bar y)
                                     \land H_i(\bar u;\bar y)\bigr),\\
 \bigwedge_{i:T_i=L}\exists\bar y
     \bigl(\bar x^i\ne\mathbf a\land S_1(\bar y,\bar x^i)
                                     \land H_i(\bar u;\bar y)\bigr).
 \end{gathered}
\]
Tuple inequality abbreviates the disjunction of coordinate inequalities.
Both tape boundaries are explicit here; a head does not move from the last
square to the first via the cyclic successor.  The right guard does not
eliminate a bounded accepting run because the exponent was chosen with slack.
The left guard matches the source's rejecting off-tape convention.
Nondeterminism is the finite disjunction of whole moves, so all tapes and the
next state must follow one common transition, not separately selected moves.

\paragraph{Why the tableau is equivalent to acceptance.}
Given an accepting bounded run, impose any order on the universe and identify
its elements with \(1,\ldots,n\).  Isomorphism closure says this relabeling
does not change membership in \(\mathcal C\).  Interpret the codes, digit
table, time, tapes, state and heads as just specified.  Fill the remaining
rows by freezing the accepting final configuration.  The slack leaves room
for the entire run, so every conjunct holds and the last row is accepting.

Conversely, take any expansion satisfying the conjuncts.  The order and digit
induction determine the correct input words.  The range constraints and unique
heads make each row an actual instantaneous description of \(M\).  Each
successor pair before the final row either freezes a final description or
satisfies \(\Phi_s\) for one legal move.  Induction on the finite order of
time tuples therefore reconstructs a genuine run.  Its last state is
\(q_A\), and a rejecting final state cannot become accepting because it is
frozen.  Thus \(M\) accepts the encoding of the ordered presentation of
the reduct.  Its language is exactly \(E(\mathcal C)\), so the reduct
belongs to \(\mathcal C\), again by isomorphism closure.
All functions in this finite auxiliary type can be replaced by the graph
translation already proved.  The result is a single ESO sentence over the
original input type, on the same universe.

\paragraph{The finitely many small sizes.}
Let \(n_0\) cover the code-size and slack exceptions.  For each
\(1\leq n<n_0\) there are finitely many \(S\)-structures on that size,
since \(S\) is finite and relational.  For a selected such structure
\(\mathcal D\) of size \(d\), its first-order complete finite diagram is
\[
 \exists x_1\cdots x_d\left(
   \bigwedge_{i<j}x_i\ne x_j\ \land
   \forall y\bigvee_i y=x_i\ \land\ \Delta_{\mathcal D}(\bar x)
 \right),
\]
where \(\Delta_{\mathcal D}\) contains each positive or negative relation
instance from \(\mathcal D\).  It defines its isomorphism class exactly.
Disjoin these diagrams for the finitely many small structures in
\(\mathcal C\); call the resulting sentence \(\chi\), taking an
unsatisfiable sentence if none qualify.  To select the diagrams effectively,
simulate the clocked machine on each of their finitely many encodings and
enumerate its finite computation tree to the clock bound.
Conjoin the large-universe tableau
with the first-order assertion of at least \(n_0\) distinct elements, and
disjoin \(\chi\).  If \(\bar R\) is its auxiliary type, then
\(\chi\lor\exists\bar R\theta\) is equivalent to
\(\exists\bar R(\chi\lor\theta)\): on a nonempty universe the relations
and total functions always have some interpretations.  This is one
generalized spectrum, not an infinite union.  In the empty type the same
diagrams reduce to exact-size sentences.  This proves finite modification,
closes every small case and completes both original clauses.

Finally, converting a unary size word to binary and back costs polynomial in
\(n\), and both displayed machine directions have polynomial-in-\(n\)
bounds.  Thus the unary reformulation is a consequence of the proof, while
the original binary empty-type clause remains the \(NP_1\) clause.
\end{proof}
\end{annexitem}

\begin{annexitem}{annex:p9.n-c-004.1.unordered-boundary}{thm:p09-fagin74-unordered-boundary}{Pair certificates and the failure of unordered fixed-point parity}
\begin{proof}
A fixed-point-free involution divides a finite universe into disjoint
two-element orbits, so its size is even.  Conversely on any positive even
size, partition the universe into pairs and interchange the members of every
pair.  The relational version has exactly one neighbour for each point,
symmetry makes this neighbour map an involution, and irreflexivity removes
one-element orbits.  It is the same witness in both directions.  Even size
is also polynomial-time decidable: read an explicitly represented universe
and count its elements modulo two.  On the original nonempty-type encoding,
parity is the low bit of the binary \(n\) prefix; on a binary numerical
input it is likewise trivial.  This does not assume that the structure itself
contains a definable order.

For nondefinability fix an order-free FO(LFP) sentence in the equality-only
language.  Rename its individual variables so distinct binders use distinct
names, and choose \(v\) at least their total number, including fixed-point
tuple variables and parameters.  A partial bijection between two pure sets
transports assignments of these variables exactly when it preserves their
equality pattern.  For sets of sizes at least \(v\), every extension of an
assignment needed by a quantifier can be matched: an old value goes to its
matched value; a new value can be chosen outside the at most \(v-1\)
already assigned values.  The argument works in both directions.

We prove formula invariance under these matching assignments, allowing
relation variables whose membership already depends only on the equality
pattern of their arguments and named parameters, with the same pattern
table in the two sets.  Equality atoms and such relation atoms agree;
negation and the Boolean connectives preserve agreement.  The preceding
extension argument proves both first-order quantifier cases.

For a least-fixed-point subformula, fix matching assignments to its
parameters.  Start its defined relation empty in both sets.  By induction on
the stage number, the preceding formula invariance applied to the defining
formula says that the next stage has the same membership table of equality
patterns of tuple entries and parameters.  Permutations fixing the parameters
show at each stage that every tuple of a given pattern has the same membership.
Each pattern realizable in one set is realizable in the other, since the
tuple and parameter slots together use at most \(v\) values.  There are
only finitely many such patterns.  Positivity makes the operator monotone,
so iteration from the empty relation is increasing and stabilizes after at
most that many pattern additions.  Both sets have the same finite transition
table on patterns, hence the same stabilized relation.  This proves the
fixed-point case as well.  Nested fixed points are covered by structural
induction on subformulas, not by assuming their result at the outset.

There are no free parameters for the sentence.  It consequently has one
constant truth value on all pure sets of sizes at least \(v\).  Infinitely
many such sizes are even and infinitely many are odd.  No such sentence
defines parity.  With an arbitrary fixed relational input type, interpret
every input relation as empty.  Its atomic formulas are uniformly false,
and the same induction applies.  If parity were definable on all structures
of that type it would be definable on this subclass, a contradiction.
\end{proof}
\end{annexitem}
